\section{Graph consequences of recursive upper-array norms}
\label{met:f:applications}

Recursive upper arrays give three further target norms. The first uses
a height-dependent child exponent and an outer Hilbert sum; its sixth-
and cubic-power estimates yield additive losses in squared slope norm.
The other two use Euclidean child aggregation on a joined tree and
differ at the joining root. We retain the qualitative propagation
argument for the free-root norm and an explicit finite-error cubic
argument for the fixed-zero-root norm. The constants below belong to
these stated inputs and deductions.

Let $T_h=\bigcup_{j=0}^h\mathbb N^j$ and $p_h=1+1/h$.
For real $v\in\ell_{p_h}(T_h)$, let $P(v),Q(v)$ be the least nonnegative
arrays whose difference is $v$ and which satisfy
\[
 P_s(v)\ge\Big(\sum_{t\in C(s)}P_t(v)^{p_h}\Big)^{1/p_h},
 \qquad
 Q_s(v)\ge\Big(\sum_{t\in C(s)}Q_t(v)^{p_h}\Big)^{1/p_h}.
\]
Here $C(s)$ is the set of children of $s$ and the sums at leaves are
zero.  More explicitly, if the two child norms are $a_s,b_s$, then
$P_s(v)=\max\{a_s,v_s+b_s\}$ and
$Q_s(v)=\max\{a_s-v_s,b_s\}$.  Write
\[
 \nu_h(v)=P_\emptyset(v)+Q_\emptyset(v),\qquad
 X_{\mathrm v}=\Big(\bigoplus_{h\ge1}(\ell_{p_h}(T_h),\nu_h)
                  \Big)_{\ell_2}.
\]
The intrinsic upper-array results supply completeness, reflexivity,
separability of $X_{\mathrm v}^*$, density of finite coordinate vectors,
and the estimates used below; see \cite[Section~7]{RecursivePotentials}.
We use the slopes from~\eqref{met:slopes}. To discuss paired children,
write
\begin{equation}\label{met:f:slopes}
 y_j=x+z_j,\qquad \bar y_j=x-z_j.
\end{equation}
In noncontracting normalization the $y_j$, and also the $z_j$, are
$2$-separated. In upper-one normalization the separation is $2/D$.

\subsection{The sixth-power estimate}

A finite head in $X_{\mathrm v}$ means a finite ancestor-closed set in
each of finitely many components, containing the component root whenever
that component is used.  Its coordinate projection is a bounded
finite-rank operator.  No estimate on its operator norm is needed in
the following argument.

The finite-head estimate in \cite[Theorem~7.4]{RecursivePotentials}
allows arbitrary completed tails.

\begin{proposition}\label{met:f:sixth}
For $X_{\mathrm v}$ one has
\begin{equation}\label{met:f:sixth-input}
 S-2\|x\|_{\mathrm v}^2
 \ge c_6S^{-2}\|y\|_{\mathrm v}^6,
 \qquad
 S=\|x+y\|_{\mathrm v}^2+\|x-y\|_{\mathrm v}^2,
 \qquad c_6=2\exp(-(\mathrm e-1))4^{-6},
\end{equation}
whenever $x$ is supported on a finite head and $y$ vanishes there.
The right side is interpreted as zero when $S=0$.
Every distortion-$D$ embedding of $\mathcal D_k$ into this norm satisfies
\begin{equation}\label{met:f:sixth-power-depth}
 \frac{c_6 k}{8}\le D^6-D^4.
\end{equation}
For $0<\varepsilon\le1$, put
\begin{equation}\label{met:f:sixth-q}
 q_6(\varepsilon)=
 \left(1-\frac{c_6\varepsilon^6}{512}\right)^{1/2}.
\end{equation}
Every infinite family satisfying
\[
 \|x+u_j\|_{\mathrm v},\ \|x-u_j\|_{\mathrm v}\le1,
 \qquad \|u_j-u_l\|_{\mathrm v}\ge\varepsilon\quad(j\ne l)
\]
has $\|x\|_{\mathrm v}\le q_6(\varepsilon)$.
Keeping a fixed unit-lens contraction gives the additional form
\begin{equation}\label{met:f:sixth-depth}
 D^{-1}\le q_6(D^{-1})^k.
\end{equation}
\end{proposition}

\begin{proof}
If $\|x\pm y\|_{\mathrm v}\le R$ with $R>0$, then $S\le2R^2$.
For $S>0$, \eqref{met:f:sixth-input} implies
\[
 R^2\ge\|x\|_{\mathrm v}^2+
             \frac{c_6}{8R^4}\|y\|_{\mathrm v}^6.
\]
It is trivial when $S=0$. Theorem~\ref{met:power-transfer}, with
power $6$ and coefficient $c_6/8$, gives
\eqref{met:f:sixth-power-depth}.

For the unit-lens statement, Lemma~\ref{met:finite-pair} supplies
pure head--tail pairs $x',y'$ with
$\|x'-x\|_{\mathrm v}<\tau$,
$\|y'\|_{\mathrm v}\ge\varepsilon/2-\tau$, and endpoints of
norm at most $1+2\tau$. Apply the displayed homogeneous estimate
and let $\tau\downarrow0$. It gives
$\|x\|_{\mathrm v}^2+c_6\varepsilon^6/512\le1$.
In upper-one graph normalization use the fixed separation parameter
$\varepsilon=1/D$; actual sibling separation is at least $2/D$.
Proposition~\ref{met:backward} proves \eqref{met:f:sixth-depth}.
\end{proof}

The power bound keeps the child radius in the loss at every level.
The geometric form instead uses one contraction factor chosen at the
start from $D$.

\subsection{The cubic estimate and weakly null paired differences}

The intrinsic cubic inequality and its weakly null consequence are
\cite[Theorem~7.6 and Proposition~7.7]{RecursivePotentials}.
The first gives another additive power bound. The second exhibits a
different selection: pairing a weakly Cauchy subsequence. The paired
differences and crossed convex averages also occur in the approximate
midpoint argument of \cite[proof of Corollary~2.3]{DKRRZ2016}; here
dual separability supplies the subsequence by diagonal extraction.

\begin{proposition}\label{met:f:cubic-weak}
The dual $X_{\mathrm v}^*$ is separable, and under the same
finite-head support conditions as above one has
\begin{equation}\label{met:f:cubic-input}
 \|x\pm y\|_{\mathrm v}\le R
 \quad\Longrightarrow\quad
 R^2\ge\|x\|_{\mathrm v}^2+
       \frac{c_3}{2\sqrt2\,R}\|y\|_{\mathrm v}^3,
 \qquad R>0,\quad c_3=(6\mathrm e)^{-3}.
\end{equation}
Consequently every distortion-$D$ embedding of $\mathcal D_k$ satisfies
\begin{equation}\label{met:f:cubic-power-depth}
 \frac{c_3 k}{2\sqrt2}\le D^3-D.
\end{equation}
For $0<\varepsilon\le1$, define
\begin{equation}\label{met:f:cubic-theta}
 \theta(\varepsilon)=
 \left(1-\frac{c_3\varepsilon^3}{2\sqrt2}\right)^{1/2}.
\end{equation}
If $(h_j)$ is weakly null, $\|h_j\|_{\mathrm v}\ge\varepsilon$,
and $\|x\pm h_j\|_{\mathrm v}\le1$, then
$\|x\|_{\mathrm v}\le\theta(\varepsilon)$.
Every distortion-$D$ embedding of the depth-$k$ countably branching
diamond into $X_{\mathrm v}$ satisfies
\begin{equation}\label{met:f:cubic-depth}
 1\le D\theta(D^{-1})^k.
\end{equation}
\end{proposition}

\begin{proof}
Theorem~\ref{met:power-transfer}, applied directly to
\eqref{met:f:cubic-input}, gives \eqref{met:f:cubic-power-depth}.
The weakly null implication is the cited Proposition~7.7.
To use it in the graph, normalize to lower factor one and upper
factor $D$, and let $R\le D$ bound the children of one parent $x$.
Dual separability gives a weakly Cauchy subsequence of $(y_j)$.
The diagonal argument is the one in Section~\ref{met:weak-linear}.

Take successive disjoint pairs $y_{j_{2r}},y_{j_{2r+1}}$ from this
subsequence and put
\[
 h_r=\frac{y_{j_{2r}}-y_{j_{2r+1}}}{2R}.
\]
The sequence $(h_r)$ is weakly null.  Sibling separation gives
$\|h_r\|_{\mathrm v}\ge1/R\ge1/D$.  The identities
\[
 \frac{x}{R}+h_r
   =\frac{y_{j_{2r}}+\bar y_{j_{2r+1}}}{2R},
 \qquad
 \frac{x}{R}-h_r
   =\frac{y_{j_{2r+1}}+\bar y_{j_{2r}}}{2R}
\]
show that both points belong to the unit ball. The weakly null
conclusion, with $\varepsilon=1/D$, gives $\|x\|_{\mathrm v}
\le R\theta(D^{-1})$. Taking the parent suprema and iterating
this contraction proves \eqref{met:f:cubic-depth}.
\end{proof}

\subsection{A qualitative contraction for the free-root sum norm}

We next change the norm.  Join disjoint copies of $T_h$ at a new root
$\rho$ whose children are their roots.  At every node use the least
nonnegative arrays $p,q$ with difference equal to the coordinate and
with each array dominating the $\ell_2$ norm of its child values.
For a finitely supported vector with an unrestricted root coordinate,
put
\[
 N_J(x)=p_\rho(x)+q_\rho(x),
\]
and let $X_J$ be the completion.  In this subsection
$X_J$ refers to the tree formed by joining these finite-height components,
with a free root coordinate. In particular, the norm of an indicator of an initial
path starting at a component root is two.

Restriction $P_H$ to a finite ancestral set containing $\rho$ is
contractive. On finitely supported arrays, restriction removes child
values outside $H$; the recursion is monotone in their nonnegative
norms, so neither root field can increase. This proves contractivity
on the dense finite arrays and hence on the completion. Applying
$I-P_H$ to finite approximations also proves density of finitely
supported vectors in $\ker P_H$.

The intrinsic probabilistic propagation estimate
\cite[Lemmas~6.1 and~6.2]{RecursivePotentials} supplies, for every
$0<\zeta\le2$, a number $\gamma(2,\zeta/4)>0$ such that
\begin{equation}\label{met:f:sum-input}
 \frac{N_J(x+y)+N_J(x-y)}2
 \ge N_J(x)+2\gamma(2,\zeta/4)
\end{equation}
whenever $x,y$ are finitely supported, $x$ is supported on a finite
ancestor-closed set $H$ containing $\rho$, $y$ vanishes on $H$,
$N_J(x\pm y)\le2$, and $N_J(y)\ge\zeta$.
The number $\gamma$ is uniform in the finite tree and its height.
The estimate comes from the logarithmic-martingale propagation of the
two upper arrays; only its positivity is used here.

The resulting separated-lens statement is also
\cite[Theorem~6.3]{RecursivePotentials}.

\begin{proposition}\label{met:f:sum}
For every $0<\varepsilon\le2$ there is
$\eta(\varepsilon)\in(0,1)$ such that
\[
 N_J(x\pm z_j)\le1,
 \qquad N_J(z_j-z_l)\ge\varepsilon\quad(j\ne l)
 \quad\Longrightarrow\quad
 N_J(x)\le1-\eta(\varepsilon).
\]
For $\varepsilon>2$ the hypotheses are impossible, so any
$\eta(\varepsilon)\in(0,1)$ is admissible.
An embedding of the depth-$k$ countably branching diamond into
$X_J$ with distortion at most $D$ therefore satisfies
\begin{equation}\label{met:f:sum-depth}
 D^{-1}\le\bigl(1-\eta(2/D)\bigr)^k.
\end{equation}
\end{proposition}

\begin{proof}
Fix $\varepsilon$ and put $\zeta=\varepsilon/4$.
For each sufficiently small $\delta>0$, take a finite coordinate
vector $x'$ with $N_J(x-x')<\delta$ and a finite
ancestor-closed head $H$ containing its support and $\rho$.
The sequence $(z_j)$ is bounded.  Finite-dimensional total boundedness
therefore supplies a pair with
\[
 u=(z_j-z_l)/2,
 \qquad N_J(u)\ge\varepsilon/2,
 \qquad N_J(P_Hu)<\delta.
\]
Also $N_J(x\pm u)\le1$ by convexity.  Approximate
$(I-P_H)u$ by a finitely supported vector $y$ vanishing on $H$ so
that $N_J(y-u)<2\delta$.  This is possible by first using
density and then applying the bounded operator $I-P_H$; no
contractivity is assumed.  For small enough $\delta$ we have
\[
 N_J(y)\ge\zeta,
 \qquad N_J(x'\pm y)\le1+3\delta\le2.
\]
Now \eqref{met:f:sum-input} gives
\[
 1+3\delta\ge N_J(x')+2\gamma(2,\varepsilon/16)
 \ge N_J(x)-\delta+2\gamma(2,\varepsilon/16).
\]
Let $\delta\downarrow0$.  For example,
$\eta(\varepsilon)=\min\{\gamma(2,\varepsilon/16),1/2\}$
has the required property.  This choice records positivity, not a
numerical power law in $\varepsilon$.

In upper-one normalization use the fixed separation $2/D$ and
apply Proposition~\ref{met:backward} with
$\rho=1-\eta(2/D)$. This proves \eqref{met:f:sum-depth}.
\end{proof}

\subsection{An additive cubic loss for the fixed-zero-root norm}

On the same joined tree, now require $x_\rho=0$ and define
\[
 N_0(x)=p_\rho(x)=q_\rho(x)
 =\max\left\{
       \Big(\sum_h p_{r_h}(x)^2\Big)^{1/2},
       \Big(\sum_h q_{r_h}(x)^2\Big)^{1/2}\right\},
\]
where $r_h$ is the root of the $h$th component.  Denote its
completion by $X_0$.  The root coordinate is fixed at zero and the
cost is the common root value.  On the root-zero subspace the
preceding sum norm is $2N_0$; keeping the actual normalization is
essential for the constants below.
The same ancestral projections preserve this subspace and are
contractive for $N_0$; finitely supported tails are dense here as well.

The deterministic upper-array estimate for $X_0$ is
\cite[Theorem~2.1, equation~(2.4)]{RecursivePotentials}:
\begin{equation}\label{met:f:zero-input}
 \frac{N_0(x+y)+N_0(x-y)}2
 \ge N_0(x)+\frac{N_0(y)^3}{8(2N_0(x)+N_0(y))^2},
\end{equation}
for every finite ancestor-closed head containing $\rho$, every $x$
supported on it, and every nonzero completed tail $y$ vanishing
there.  Finite support is required of $x$, but not of $y$.

\begin{proposition}\label{met:f:zero}
If the depth-$k$ countably branching diamond admits an embedding into
$X_0$ with distortion at most $D\ge1$, then
\begin{equation}\label{met:f:zero-depth}
 k\le4608D^3(1-D^{-1}).
\end{equation}
\end{proposition}

\begin{proof}
Normalize to lower constant $A=1/D$ and upper constant one.
For each level $i$, let $L_i$ be the supremum of the $N_0$-norms
of all its historical edge slopes.  Thus $A\le L_i\le1$.
Fix any parent slope $x$ at level $i<k$.  Its child slopes satisfy
$N_0(y_j),N_0(\bar y_j)\le L_{i+1}$ and
$N_0(y_j-y_l)\ge2A$ for $j\ne l$.
For such a pair put $z=(y_j-y_l)/2$.  Then
\begin{equation}\label{met:f:zero-pair}
 A\le N_0(z)\le1,
 \qquad N_0(x\pm z)\le L_{i+1},
\end{equation}
because $x+z=(y_j+\bar y_l)/2$ and
$x-z=(y_l+\bar y_j)/2$.

Fix the errors before selecting the finite head and the pair:
\[
 c_0=\frac{A^3}{2304},\qquad
 \eta=\frac{c_0}{6}.
\]
Since $0<A\le1$, we have $\eta\le A/2$ and $\eta\le1$.
Choose a finitely supported $x'$ with $N_0(x-x')\le\eta$, and
then a finite ancestor-closed head $H$ containing its support and
$\rho$.  The projected child slopes $P_Hy_j$ form a bounded subset
of a finite-dimensional space.  Choose $j\ne l$ so that their
half-difference $z$ satisfies $N_0(P_Hz)\le\eta$, and put
$y=z-P_Hz$.  Then
\[
 A/2\le N_0(y)\le1+\eta,
 \qquad N_0(x')\le1+\eta.
\]
In particular $y\ne0$, so \eqref{met:f:zero-input} applies to
$x',y$.  Its gain is at least
\[
 \frac{(A/2)^3}{8\,[3(1+\eta)]^2}
 \ge\frac{(A/2)^3}{8\cdot6^2}
 =\frac{A^3}{2304}=c_0.
\]
For both signs,
$N_0(x\pm z)\ge N_0(x'\pm y)-2\eta$,
and $N_0(x')\ge N_0(x)-\eta$.  Consequently
\[
 \begin{split}
 L_{i+1}
 &\ge\frac{N_0(x+z)+N_0(x-z)}2\\
 &\ge N_0(x)+c_0-3\eta
  =N_0(x)+\frac{c_0}{2}.
 \end{split}
\]
The parent edge was arbitrary.  Taking the supremum over all parent
edges gives $L_i\le L_{i+1}-c_0/2$, without requiring that a
supremum be attained.  Iteration yields
\[
 A\le L_0\le L_k-\frac{kc_0}{2}
 \le1-\frac{kA^3}{4608},
\]
which is \eqref{met:f:zero-depth}.
\end{proof}

The bounds for $X_{\mathrm v}$ distinguish an additive power budget
from a contraction chosen at a fixed separation. For the Euclidean
joined-tree models, \eqref{met:f:sum-depth} records the positivity
obtained by probabilistic propagation, while
\eqref{met:f:zero-depth} retains the explicit approximation errors
in a deterministic cubic-gain argument. Each conclusion rules out a
distortion bound independent of the diamond depth, in its stated norm.
